\documentclass[10pt,a4paper]{amsart}
\usepackage{amsmath,amssymb,amsthm,graphicx}
\usepackage{verbatim}
\newenvironment{reference}{\comment}{\endcomment}
%\newenvironment{reference}{}{}
\newenvironment{slogan}{\comment}{\endcomment}
\newenvironment{history}{\comment}{\endcomment}

% For commutative diagrams we use Xy-pic
\usepackage[all]{xy}

% We use 2cell for 2-commutative diagrams.
\xyoption{2cell}
% Theorem environments.
%
\theoremstyle{plain}
\newtheorem{theorem}[subsection]{Theorem}
\newtheorem{proposition}[subsection]{Proposition}
\newtheorem{lemma}[subsection]{Lemma}

\theoremstyle{definition}
\newtheorem{definition}[subsection]{Definition}
\newtheorem{example}[subsection]{Example}
\newtheorem{exercise}[subsection]{Exercise}
\newtheorem{situation}[subsection]{Situation}

\theoremstyle{remark}
\newtheorem{remark}[subsection]{Remark}
\newtheorem{remarks}[subsection]{Remarks}

\numberwithin{equation}{subsection}

% Macros
%
\def\lim{\mathop{\mathrm{lim}}\nolimits}
\def\colim{\mathop{\mathrm{colim}}\nolimits}
\def\Spec{\mathop{\mathrm{Spec}}}
\def\Hom{\mathop{\mathrm{Hom}}\nolimits}
\def\Ext{\mathop{\mathrm{Ext}}\nolimits}
\def\SheafHom{\mathop{\mathcal{H}\!\mathit{om}}\nolimits}
\def\SheafExt{\mathop{\mathcal{E}\!\mathit{xt}}\nolimits}
\def\Sch{\mathit{Sch}}
\def\Mor{\mathop{\mathrm{Mor}}\nolimits}
\def\Ob{\mathop{\mathrm{Ob}}\nolimits}
\def\Sh{\mathop{\mathit{Sh}}\nolimits}
\def\NL{\mathop{N\!L}\nolimits}
\def\CH{\mathop{\mathrm{CH}}\nolimits}
\def\proetale{{pro\text{-}\acute{e}tale}}
\def\etale{{\acute{e}tale}}
\def\QCoh{\mathit{QCoh}}
\def\Ker{\mathop{\mathrm{Ker}}}
\def\Im{\mathop{\mathrm{Im}}}
\def\Coker{\mathop{\mathrm{Coker}}}
\def\Coim{\mathop{\mathrm{Coim}}}

% Boxtimes
%
\DeclareMathSymbol{\boxtimes}{\mathbin}{AMSa}{"02}

%
% Macros for moduli stacks/spaces
%
\def\QCohstack{\mathcal{QC}\!\mathit{oh}}
\def\Cohstack{\mathcal{C}\!\mathit{oh}}
\def\Spacesstack{\mathcal{S}\!\mathit{paces}}
\def\Quotfunctor{\mathrm{Quot}}
\def\Hilbfunctor{\mathrm{Hilb}}
\def\Curvesstack{\mathcal{C}\!\mathit{urves}}
\def\Polarizedstack{\mathcal{P}\!\mathit{olarized}}
\def\Complexesstack{\mathcal{C}\!\mathit{omplexes}}
% \Pic is the operator that assigns to X its picard group, usage \Pic(X)
% \Picardstack_{X/B} denotes the Picard stack of X over B
% \Picardfunctor_{X/B} denotes the Picard functor of X over B
\def\Pic{\mathop{\mathrm{Pic}}\nolimits}
\def\Picardstack{\mathcal{P}\!\mathit{ic}}
\def\Picardfunctor{\mathrm{Pic}}
\def\Deformationcategory{\mathcal{D}\!\mathit{ef}}
\usepackage{cleveref}
\crefname{originaltheorem}{original theorem}{original theorems}
\Crefname{originaltheorem}{Original theorem}{Original theorems}
\makeatletter
\newlabel{bootstrap-theorem-final-bootstrap}{{04S6}{source}{}{sourcepointer-1}{}}
\newlabel{bootstrap-theorem-final-bootstrap@cref}{{[originaltheorem][0][]04S6}{[1][0][]source}{}{}{}}
\makeatother
\crefname{originalsection}{original section}{original sections}
\Crefname{originalsection}{Original section}{Original sections}
\makeatletter
\newlabel{formal-spaces-section-completion}{{0AIX}{source}{}{sourcepointer-2}{}}
\newlabel{formal-spaces-section-completion@cref}{{[originalsection][0][]0AIX}{[1][0][]source}{}{}{}}
\makeatother
\crefname{originalsection}{original section}{original sections}
\Crefname{originalsection}{Original section}{Original sections}
\makeatletter
\newlabel{section-completion-functor}{{0GDB}{source}{}{sourcepointer-3}{}}
\newlabel{section-completion-functor@cref}{{[originalsection][0][]0GDB}{[1][0][]source}{}{}{}}
\makeatother
\crefname{originallemma}{original lemma}{original lemmas}
\Crefname{originallemma}{Original lemma}{Original lemmas}
\makeatletter
\newlabel{spaces-lemma-glueing-algebraic-spaces}{{02WR}{source}{}{sourcepointer-4}{}}
\newlabel{spaces-lemma-glueing-algebraic-spaces@cref}{{[originallemma][0][]02WR}{[1][0][]source}{}{}{}}
\makeatother
\title{The smash construction factors a morphism through one etale outside T and isomorphic on completions}
\author{The Stacks Project Authors}
\date{stacks-project:restricted.tex:local-snapshot-20261012; GNU FDL 1.2 or later}
\begin{document}
\maketitle
\noindent\textit{Editorial note.} All delivered mathematical text is the author's own, copied from the pinned Stacks Project chapter file restricted.tex; no statement or proof step was rewritten, reordered or removed.\par\medskip
\noindent\textit{Editorial note.} The counted proof is the author's complete argument: the source proves the lemma in full, including the etale equivalence relation check and the completion isomorphism\par\medskip
\noindent\textit{Editorial note.} The label lemma-smash-away-from-T names the counted statement, and the proof environment immediately following it in the source is delivered verbatim as its proof.\par\medskip
\noindent\textit{Editorial note.} The counted statement is carried by the author's own lemma environment in the chapter source.\par\medskip
\noindent\textit{Allowed standard prerequisites (editorial).}\begin{itemize}
\item formal completion along a closed subset and the completion functor of the chapter
\item glueing of algebraic spaces along open immersions to produce pushout spaces
\item etale equivalence relations and the bootstrap quotient theorem
\item checking finite type and isomorphism conditions etale locally
\end{itemize}\medskip
\ifdefined\hypertarget\hypertarget{sourcepointer-1}{}\fi
\noindent\textit{Source reference.} Original theorem 04S6: sources/stacks-chow/tags.txt, source lines 6209--6209. Prerequisite pointer invoked at proof line(s) 6595; it names a standard fact of a sibling chapter that the counted argument uses, and every step of the counted proof built on it is retained verbatim in the delivered body.\par\medskip
\ifdefined\hypertarget\hypertarget{sourcepointer-2}{}\fi
\noindent\textit{Source reference.} Original section 0AIX: sources/stacks-chow/tags.txt, source lines 13280--13280. Prerequisite pointer invoked at proof line(s) 6607; it names a standard fact of a sibling chapter that the counted argument uses, and every step of the counted proof built on it is retained verbatim in the delivered body.\par\medskip
\ifdefined\hypertarget\hypertarget{sourcepointer-3}{}\fi
\noindent\textit{Source reference.} Original section 0GDB: sources/stacks-restricted/pinned-original-source.tex, source lines 5953--5953. Prerequisite pointer invoked at proof line(s) 6572; it names a standard fact of the same chapter that the counted argument uses, and every step of the counted proof built on it is retained verbatim in the delivered body.\par\medskip
\ifdefined\hypertarget\hypertarget{sourcepointer-4}{}\fi
\noindent\textit{Source reference.} Original lemma 02WR: sources/stacks-chow/tags.txt, source lines 3814--3814. Prerequisite pointer invoked at proof line(s) 6589; it names a standard fact of a sibling chapter that the counted argument uses, and every step of the counted proof built on it is retained verbatim in the delivered body.\par\medskip
\begin{lemma}
\label{lemma-smash-away-from-T}
Let $S$ be a scheme. Let $X$ be a locally Noetherian algebraic space over $S$
and let $T \subset |X|$ be a closed subset. Let $f : X' \to X$ be a morphism
of algebraic spaces which is locally of finite type and \'etale outside of $T$.
There exists a factorization
$$
X' \longrightarrow X'' \longrightarrow X
$$
of $f$ with the following properties:
$X'' \to X$ is locally of finite type,
$X'' \to X$ is an isomorphism over $X \setminus T$, and
$X'_{/T} \to X''_{/T}$ is an isomorphism (see proof for notation).
\end{lemma}

\begin{proof}
The notation using the subscript ${}_{/T}$ in the statement refers to the
construction which to a morphism $f : X' \to X$ of algebraic spaces
associates the morphism $f_{/T} : X'_{/f^{-1}T} \to X_{/T}$ of formal
algebraic spaces, see Section \ref{section-completion-functor}.
We will also use the notion $U \subset X$ and $U' \subset X'$ to denote
the open subspaces with $|U| = |X| \setminus T$ and
$U' = |X'| \setminus f^{-1}T$ introduced in
Section \ref{section-completion-functor}.

\medskip\noindent
After replacing $X'$ by $X' \amalg U$ we may and do assume
the image of $X' \to X$ contains $U$.
Let
$$
R = X' \amalg_{U'} (U' \times_U U')
$$
be the pushout of $U' \to X'$ and the diagonal morphism
$U' \to U' \times_U U' = U' \times_X U'$. Since $U' \to X$ is \'etale,
this diagonal is an open immersion and we see that $R$ is an algebraic space
(this follows for example from
Spaces, Lemma \ref{spaces-lemma-glueing-algebraic-spaces}).
The two projections $U' \times_U U' \to U'$ extend to $R$
and we obtain two \'etale morphisms $s, t : R \to X'$.
Checking on each piece separately we find that $R$
is an \'etale equivalence relation on $X'$. Set $X'' = X'/R$
which is an algebraic space by
Bootstrap, Theorem \ref{bootstrap-theorem-final-bootstrap}.
By construction have the factorization as in the lemma and
the morphism $X'' \to X$ is locally of finite type (as this
can be checked \'etale locally, i.e., on $X'$).
Since $U' \to U$ is a surjective \'etale morphism
and since $s^{-1}(U') = t^{-1}(U') = U' \times_U U'$
we see that $U'' = U \times_X X'' \to U$ is an isomorphism.
Finally, we have to show the morphism $X' \to X''$ induces an isomorphism
$X'_{/T} \to X''_{/T}$. To see this, note that the formal completion of $R$
along the inverse image of $T$ is equal to the formal completion of
$X'$ along the inverse image of $T$ by our choice of $R$! By
our construction of the formal completion in
Formal Spaces, Section \ref{formal-spaces-section-completion}
we have $X''_{/T} = (X'_{/T}) / (R_{/T})$ as sheaves. Since
$X'_{/T} = R_{/T}$ we conclude that $X'_{/T} = X''_{/T}$
and this finishes the proof.
\end{proof}

\end{document}
